\section{Describing a sparse incidence pair}
\label{sec:predictor-statement}

The following theorem supplies the geometric input to compression.
An encoder knows two hidden sets. A decoder knows larger sets
containing them and a collection of independent public random
tables. A message describes a set the decoder can reconstruct.
No bound on the time needed to search or decode is asserted.

For precision, a public table is a sequence of independent proposed
sample rows, with its distribution specified by the decoder's
current information. Uniform samples in any finite known set may
be generated by independent uniforms and a fixed ordering of that
set. Distinct calls use independent tables. The encoder sends a
row index and any specified additional data; the decoder reproduces
that row without knowing either hidden set.

\begin{theorem}[Sparse-pair description]
\label{thm:predictor}
Fix \(0<\eta<1/10\), set \(\beta=\eta/10^7\), and let \(q\) be a
sufficiently large prime. Write \(\sigma=\log q\), and suppose
\[
 \sigma^\beta\le D\le\sigma^{1-\eta/2},\qquad
 L=D\sigma^{8\beta},\qquad
 R\in[\sigma^\beta,2\sigma^\beta+2]\cap2\mathbb N,\qquad P=LR.
\]
Let \(2\le d\le4\), and let
\(S\subseteq U_S\subseteq\mathcal P_d\) and
\(T\subseteq U_T\subseteq\mathcal P_d^*\) be nonempty. Assume
\[
 |S||T|\ge q^{d+1}e^{-b},\qquad
 0\le b\le C_bD\sigma^{6\beta},\qquad
 \frac{e(S,T)}{|S||T|}\le\frac{\tau}{q},\qquad
 0<\tau\le\sigma^{-200\cdot2^{d-2}\beta},
\]
where \(C_b\) is a fixed constant.
Put \(d_S=\log(|U_S|/|S|)\) and \(d_T=\log(|U_T|/|T|)\).
Orient the pair so that \(|S|\le|T|\).
There is a public-table scheme for which, except with probability
\(O(e^{-cq})\), the encoder can send a message of length at most
\[
 CqP(d_S+d_T+P)
\]
bits describing a set \(W\subseteq U_S\) such that
\[
 |W|\le |S|e^{CP},\qquad |W\cap S|\ge c|S|.
\]
The constants may depend on \(\eta,C_b\), but not on the hidden sets,
\(D\), or \(q\). The decoder uses only \(U_S,U_T\), the public
parameters and tables, and the message. There is no uniformity
assertion for the captured subset \(W\cap S\).
\end{theorem}

The orientation requires one bit and may be reversed during the
proof. Message lengths are measured in bits, whereas entropies are
in natural units; this only changes constants. Integers can be
sent with a self-delimiting code of length \(O(1+\log(m+1))\).
All searches below have explicit finite cutoffs.

We first use Theorem~\ref{thm:predictor} to obtain a contradiction
to the selected-tuple entropy bound. Its independent proof starts
in Section~\ref{sec:predictor-setup}.
